\section{Related work}\label{sec:related-work}

The background language of this paper is standard finite combinatorics. The type-\(A\) and symmetric-group setting is part of the Coxeter-group framework treated by Bj{\"o}rner and Brenti \cite{BjornerBrenti2005}. The poset and linear-extension terminology follows the classical treatments of Stanley and Trotter \cite{Stanley2012EC1,Trotter1992}. Stanley's order-polytope paper \cite{Stanley1986} is relevant background for the broader setting in which linear extensions index simplex pieces of order-polytopal objects, although no order-polytope theorem is proved here.

Counting and generating linear extensions are well-studied problems; see Brightwell and Winkler \cite{BrightwellWinkler1991} and Pruesse and Ruskey \cite{PruesseRuskey1994}. The finite rows used here are small enough to be checked directly, but these references place the common-precedence criterion of \cref{prop:common-precedence-exactness} in the standard linear-extension setting rather than as an isolated convention.

There is also nearby literature on quotients, splittings, and sets of permutations. Bj{\"o}rner and Wachs study generalized quotients in Coxeter groups \cite{BjornerWachs1988}. Martin and Sagan introduce a notion of transitivity for groups and sets of permutations \cite{MartinSagan2006}. These topics are adjacent to subgroup rows and permutation sets, but the claims in this paper are narrower: the fixed \(D_4\) and \(V_4\) rows are used as explicit finite witnesses, not as evidence for a general transitivity, quotient, splitting, or tiling theorem.

This paper is also related in method to the author's type-\(A\) poset-cone and extension-tiler work \cite{BabanskyyTypeAPoset2026}. That work studies larger type-\(A\) tiling questions. The present paper instead fixes the rank-three \(S_4\) layer and separates local predicates that can be confused if an indexing bridge is treated as a structural theorem.